% =====================================================================
\section{Kết quả cấp độ 1: Visual Genome $\cap$ COCO}
\label{sec:l1}

R@$k$ là tỉ lệ truy vấn có ảnh đúng trong $k$ kết quả đầu; MRR là trung bình nghịch đảo hạng của ảnh đúng, tính trong
top-50. Mọi số của mô hình huấn luyện là trung bình $\pm$ độ lệch chuẩn qua 3 seed, lấy ở checkpoint tốt nhất theo val.
Khi so hai mô hình, nhóm dùng kiểm định Welch trên 3 seed mỗi bên.

\subsection{So sánh CLIP với CLIP + đồ thị trên test}

\begin{table}[t]
\centering
\caption{Kết quả trên test VG $\cap$ COCO (10\,696 truy vấn, pool 2\,138 ảnh). Test chấm một lần cho mỗi mô hình đã chốt.}
\label{tab:l1-test}
\footnotesize
\setlength{\tabcolsep}{3pt}
\resizebox{\columnwidth}{!}{%
\begin{tabular}{lcccc}
\toprule
Phương pháp & R@1 & R@5 & R@10 & MRR \\
\midrule
CLIP thuần & 39,32 & 66,39 & 77,32 & 52,00 \\
Ba kênh, không huấn luyện & 42,54 & 70,04 & 80,35 & 55,05 \\
Mô hình cấp 1 (\texttt{main\_r3}) & \pmstd{45,77}{0,23} & \pmstd{72,87}{0,31} & \pmstd{82,69}{0,31} & \pmstd{58,14}{0,18} \\
\textbf{Mô hình chính} & \textbf{\pmstd{47,95}{0,15}} & \textbf{\pmstd{74,79}{0,03}} & \textbf{\pmstd{84,38}{0,08}} & \textbf{\pmstd{60,15}{0,15}} \\
\bottomrule
\end{tabular}}
\tabsrc{\file{experiments/level1/test/}, \file{experiments/improve/test/}}
\end{table}

Bảng~\ref{tab:l1-test} là phép so bắt buộc của đề: CLIP thuần và CLIP + đồ thị trên cùng cách chia. Dòng ``ba kênh, không
huấn luyện'' dùng cùng công thức trộn nhưng vector của phần đồ thị là vector CLIP text gốc, không qua GAT. Mô hình chính hơn
CLIP 8,63 điểm R@1. So với mô hình cấp~1, mô hình chính hơn 2,18 điểm R@1 ($p = 0{,}0005$) và 2,00 điểm MRR
($p = 0{,}0002$). Ba seed của mô hình chính cho R@1 48,06 / 47,78 / 47,99. Từ val sang test, R@1 giảm 0,96 điểm (48,91
$\to$ 47,95), so với 0,71 điểm của mô hình cấp~1; mô hình chính được xây qua nhiều bước chọn trên val
(Bảng~\ref{tab:l1-chain}).

\subsection{Mô hình chính được xây thế nào}

\begin{table}[t]
\centering
\caption{Từng bước từ mô hình cấp~1 tới mô hình chính, R@1 và MRR trên val. $p$: Welch so với bước trước.}
\label{tab:l1-chain}
\footnotesize
\setlength{\tabcolsep}{2.5pt}
\begin{tabular}{p{3.4cm}ccc}
\toprule
Bước (đổi một thứ so với bước trước) & R@1 & MRR & Chênh R@1 ($p$) \\
\midrule
Mô hình cấp 1 & \pmstd{46,48}{0,24} & 58,85 & \\
+ Đồ thị câu FACTUAL-T5 & \pmstd{47,23}{0,36} & 59,49 & +0,76 ($\approx$0,05) \\
+ Trộn điểm thô & \pmstd{47,48}{0,10} & 59,72 & +0,25 (n.s.) \\
+ Đồ thị ảnh sạch và liền\textsuperscript{*} & \pmstd{48,38}{0,17} & 60,51 & +0,90 (0,003) \\
+ Số lượng và bảng $g$ & \textbf{\pmstd{48,91}{0,21}} & \textbf{60,86} & +0,53 (0,032) \\
\bottomrule
\multicolumn{4}{p{0.95\columnwidth}}{\scriptsize \textsuperscript{*}Gộp bốn thay đổi (gộp trùng, hình học cạnh, cạnh
kNN, node tổng) theo quyết định của nhóm; ablation bỏ cạnh ở Bảng~\ref{tab:l1-ablation-main} đo phần đóng góp của cạnh.}
\end{tabular}
\tabsrc{\file{experiments/improve/README.md}}
\end{table}

Mô hình chính được xây qua bốn bước; mỗi bước đổi một thứ và được giữ lại khi R@1 val tăng qua 3 seed
(Bảng~\ref{tab:l1-chain}). Mức tăng của bước trộn điểm thô không có ý nghĩa thống kê. Bước này vẫn được giữ vì nó cho phép
tách chính xác điểm cuối thành từng phần (Mục~\ref{sec:method-expl}) và có độ lệch giữa các seed nhỏ nhất. Các hướng thử trên val nhưng không
giữ (so khớp bộ ba tách thành phần, kênh mức cả đồ thị, gom vật cùng loại thành một node nhóm, \dots) ghi ở
\file{experiments/improve/README.md}. Ví dụ, gom vật cùng loại làm R@1 giảm 0,63 điểm ($p = 0{,}019$), nên mô hình chính chỉ
thêm số lượng vào từng node.

\subsection{Ablation: cải thiện có đến từ cấu trúc quan hệ không}

Đề yêu cầu bỏ cạnh hoặc thay bằng cạnh ngẫu nhiên để kiểm tra phần cải thiện có đến từ cấu trúc quan hệ. Nhóm làm
hai loại ablation. Loại thứ nhất huấn luyện lại từ đầu, cùng công thức, cùng 3 seed. Loại thứ hai giữ nguyên checkpoint và
chỉ phá đồ thị lúc chấm.

\begin{table}[t]
\centering
\caption{Ablation của mô hình cấp~1, mỗi dòng huấn luyện lại từ đầu. R@1.}
\label{tab:l1-ablation}
\footnotesize
\setlength{\tabcolsep}{3pt}
\begin{tabular}{p{4.3cm}cc}
\toprule
Cấu hình & Val & Test \\
\midrule
Mô hình cấp 1 & \pmstd{46,48}{0,24} & \pmstd{45,77}{0,23} \\
$a,b,c$ cố định cho mọi câu & \pmstd{46,13}{0,15} & \pmstd{45,13}{0,05} \\
Bỏ GAT (không lan truyền) & \pmstd{46,71}{0,07} & \pmstd{45,95}{0,20} \\
Tắt kênh bộ ba (GAT vẫn dùng cạnh) & \pmstd{45,99}{0,05} & \pmstd{45,07}{0,22} \\
Bỏ hẳn quan hệ & \pmstd{45,98}{0,10} & \pmstd{45,26}{0,14} \\
Câu không qua bộ mã hoá đồ thị & \pmstd{46,92}{0,07} & \pmstd{45,61}{0,29} \\
Nối lại cạnh ngẫu nhiên (giữ bậc) & \pmstd{46,24}{0,24} & \pmstd{45,50}{0,40} \\
Loss LambdaRank thay cross-entropy & \pmstd{45,18}{0,34} & \pmstd{44,28}{0,30} \\
\bottomrule
\end{tabular}
\tabsrc{\file{experiments/level1/README.md}}
\end{table}

\begin{table}[t]
\centering
\caption{Ablation và đối chứng của mô hình chính, chỉ trên val. R@1 (chênh so với mô hình đủ).}
\label{tab:l1-ablation-main}
\footnotesize
\setlength{\tabcolsep}{3pt}
\begin{tabular}{p{3.5cm}cc}
\toprule
 & Mô hình chính & Mô hình cấp 1 \\
\midrule
Đủ đồ thị & \pmstd{48,91}{0,21} & \pmstd{46,48}{0,24} \\
\multicolumn{3}{l}{\emph{Huấn luyện lại}} \\
Bỏ mọi cạnh & $-$2,39 ($p$ = 0,0024) & $-$0,50 ($p$ = 0,055) \\
\multicolumn{3}{l}{\emph{Chỉ phá lúc chấm}} \\
Đồ thị của một ảnh khác & 21,91 ($-$27,00) & 26,33 ($-$20,15) \\
Nối lại quan hệ có nhãn ngẫu nhiên & 48,32 ($-$0,59) & 46,29 ($-$0,19) \\
Bỏ bảng $g$ & $-$0,35 & --- \\
\bottomrule
\end{tabular}
\tabsrc{\file{experiments/improve/no-edges/REPORT.md}, \file{experiments/improve/controls/README.md}}
\end{table}

Bảng~\ref{tab:l1-ablation} và~\ref{tab:l1-ablation-main} cho ba kết luận.
\begin{itemize}
\item \textbf{Mô hình dựa vào đồ thị của đúng ảnh.} Khi mỗi ảnh mang scene graph của một ảnh khác, R@1 val của mô hình chính
giảm từ 48,91 xuống 21,91, thấp hơn CLIP thuần (39,13).
\item \textbf{Ở mô hình cấp~1, cạnh đóng góp ít.} Bỏ hẳn quan hệ làm R@1 giảm khoảng 0,5 điểm; bỏ GAT không làm giảm kết quả.
\item \textbf{Ở mô hình chính, cạnh đóng góp 2,39 điểm.} Huấn luyện lại sau khi bỏ mọi cạnh làm R@1 val giảm 2,39 điểm
($p = 0{,}0024$), tương ứng 24\% mức chênh giữa mô hình chính và CLIP trên val (2,39 / 9,78); ở mô hình cấp~1 tỉ lệ này là
7\%. Khi bỏ cạnh, mô hình chính còn 46,52, ngang mô hình cấp~1 (46,48). Xáo các quan hệ có nhãn lúc chấm làm R@1 giảm 0,59
điểm. Phần đóng góp riêng của từng loại cạnh (quan hệ có nhãn, kNN, node tổng) chưa được tách.
\end{itemize}

\subsection{Độ nhạy theo $\alpha$ và đồ thị giúp ở đâu}

\begin{figure}[t]
\centering
\begin{tikzpicture}
\begin{axis}[width=\columnwidth, height=4.6cm, xlabel={$\alpha$ (trọng số CLIP)}, ylabel={R@1 val (\%)},
  xmin=0, xmax=1, ymin=20, ymax=52, grid=major, grid style={gray!25},
  legend style={font=\scriptsize, at={(0.02,0.03)}, anchor=south west}, tick label style={font=\scriptsize},
  label style={font=\small}]
\addplot[cObj, thick, mark=*, mark size=1.2pt] coordinates {(0,32.81) (0.1,36.77) (0.2,40.73) (0.3,44.12) (0.4,46.52)
  (0.5,47.82) (0.6,48.14) (0.7,47.37) (0.8,45.49) (0.9,42.84) (1,39.13)};
\addlegendentry{VG $\cap$ COCO, mô hình chính}
\addplot[cTri, thick, mark=square*, mark size=1.1pt] coordinates {(0,32.40) (0.05,32.40) (0.1,32.40) (0.15,32.40)
  (0.2,32.40) (0.25,32.40) (0.3,32.40) (0.35,32.40) (0.4,32.06) (0.45,31.84) (0.5,31.61) (0.55,31.50) (0.6,30.61)
  (0.65,29.93) (0.7,29.71) (0.75,29.26) (0.8,28.36) (0.85,27.91) (0.9,26.57) (0.95,25.67) (1,24.33)};
\addlegendentry{MKG-W, CLIP + đếm cạnh ($\beta=0$)}
\end{axis}
\end{tikzpicture}
\caption{R@1 val khi ép một $\alpha$ cố định cho mọi truy vấn ($\alpha=1$ là CLIP thuần). Với VG, $b$ và $c$ giữ tỉ lệ đã
học của từng câu, trung bình 3 seed. Với MKG-W, seed 0 (đường cong giống hệt ở cả 3 seed vì $\beta=0$ không dùng NativE).
Nguồn: \file{experiments/improve/alpha-groups/alpha_sweep.json}, \file{experiments/level2/val/alpha_beta_sweep.json}.}
\label{fig:alpha}
\end{figure}

\textbf{Độ nhạy theo $\alpha$} (Hình~\ref{fig:alpha}). Trên VG, đường cong có đỉnh ở $\alpha = 0{,}6$ (R@1
\pmstd{48,14}{0,30}). Trọng số học theo từng câu cho 48,91, cao hơn $\alpha$ cố định tốt nhất 0,77 điểm, dù $\alpha$ cố định
này được chọn trên chính val. Trung bình $a$ học được là \pmstd{0,57}{0,06}, gần $\alpha$ cố định tốt nhất; mức chênh 0,77
điểm đến từ việc trọng số thay đổi theo từng câu. Chỉ dùng đồ thị ($\alpha = 0$) đạt 32,81 và chỉ dùng CLIP ($\alpha = 1$) đạt
39,13, đều thấp hơn mọi giá trị trong khoảng $\alpha \in [0{,}4;\,0{,}8]$. Đường cong của MKG-W được phân tích ở
Mục~\ref{sec:l2}.

\begin{table}[t]
\centering
\caption{R@1 val theo nhóm câu. Nhóm định nghĩa một lần từ đồ thị FACTUAL của câu, dùng chung cho mọi mô hình. KTC 95\%:
bootstrap ghép cặp theo câu.}
\label{tab:l1-groups}
\footnotesize
\setlength{\tabcolsep}{2pt}
\begin{tabular}{lrcccc}
\toprule
Nhóm & Số câu & CLIP & Cấp 1 & Chính & Chính $-$ cấp 1 [KTC] \\
\midrule
Tất cả & 10\,633 & 39,13 & 46,48 & 48,91 & +2,43 [1,92; 2,93] \\
Có số lượng nêu rõ & 1\,235 & 44,21 & 50,69 & 54,06 & +3,37 [1,75; 4,91] \\
Có $\ge$1 quan hệ & 10\,342 & 39,31 & 46,83 & 49,38 & +2,55 \\
Chỉ vật thể & 291 & 32,99 & 33,91 & 32,19 & $-$1,72 [$-$4,7; 1,4] \\
2 vật thể & 3\,757 & 35,67 & 40,87 & 42,41 & +1,54 [0,72; 2,40] \\
$\ge$3 vật thể & 6\,607 & 41,40 & 50,20 & 53,32 & +3,11 [2,45; 3,77] \\
\bottomrule
\end{tabular}
\tabsrc{\file{experiments/improve/alpha-groups/REPORT.md}}
\end{table}

\textbf{Đồ thị giúp ở đâu} (Bảng~\ref{tab:l1-groups}). Đây là câu hỏi của nhiệm vụ~2 trong đề.
\begin{itemize}
\item Mức tăng so với CLIP tăng theo số vật thể trong câu: +6,74 điểm ở câu có 2 vật thể, +11,92 điểm ở câu từ 3 vật thể
trở lên.
\item Với câu chỉ nêu vật thể, không có quan hệ (2,7\% số câu), mô hình chính thấp hơn CLIP 0,80 điểm.
\item Với câu nêu số lượng, mức tăng so với mô hình cấp~1 là +3,37 điểm; với các câu còn lại là +2,31 điểm.
\item Phần mất khi bỏ cạnh cũng tập trung ở câu có quan hệ và câu nhiều vật thể; câu chỉ có vật thể gần như không đổi
(\file{experiments/improve/no-edges/REPORT.md}).
\end{itemize}

\subsection{Giải thích và phân tích 10 truy vấn}
\label{sec:l1-cases}

\textbf{Cách chọn.} Chỉ xét câu val có ít nhất một quan hệ. Nhóm bốc ngẫu nhiên với seed cố định (2026), theo quy tắc đặt
trước: 5 câu thành công (CLIP xếp ảnh đúng dưới hạng 1, mô hình xếp hạng 1), 3 câu mô hình làm tệ đi (CLIP hạng 1, mô hình
dưới hạng 1), 2 câu cả hai đều sai. Trên toàn val, mô hình chính (seed 0) có 1\,528 câu thành công, 463 câu làm tệ đi,
4\,944 câu cả hai đều sai và 3\,698 câu cả hai đều đúng. Bảng~\ref{tab:l1-cases} tóm tắt 10 câu. Ảnh, bảng giải thích đầy
đủ và lời phân tích từng câu (viết sau khi xem ảnh thật) ở \file{experiments/improve/explanations/README.md}.

\begin{table*}[t]
\centering
\caption{Mười truy vấn của mô hình chính (seed 0, val). Hạng ảnh đúng theo CLIP $\to$ theo mô hình.}
\label{tab:l1-cases}
\footnotesize
\begin{tabular}{cp{7.6cm}cp{6.0cm}}
\toprule
Ca & Truy vấn & Hạng & Nhận xét \\
\midrule
1 & Dining table for six in a wood paneled room with a fireplace. & 2 $\to$ 1 & Thành công \\
2 & A cow standing along side of a large river. & 8 $\to$ 1 & Thành công \\
3 & A man sitting in front of the pen with a dog herding sheep in the background. & 3 $\to$ 1 & Thành công \\
4 & Someone holding a doughnut that has a bite taken out of it. & 4 $\to$ 1 & Thành công \\
5 & People huddle under several umbrellas on a sidewalk area and a white dog \dots & 2 $\to$ 1 & Thành công; $g$ dùng số lượng của \emph{umbrella} và \emph{people} \\
6 & A bunch of bananas are sitting on a wooden plate. & 1 $\to$ 2 & VG gọi đĩa gỗ là \emph{tray}; thuộc tính ``wooden'' chưa được dùng \\
7 & A train riding along the tracks along a bank of dirt. & 1 $\to$ 5 & Đồ thị VG của ảnh đúng rỗng (0 node) \\
8 & A bus that is reading ``89 Portola Plaza'' is traveling on. & 1 $\to$ 3 & Chi tiết quyết định là chữ trên biển; scene graph không biểu diễn được \\
9 & Two men standing in a living room playing video game with remotes. & 13 $\to$ 17 & Số lượng khớp đúng ($g$ +0,071 cho ``two men'') nhưng ảnh đúng thiếu nhãn tay cầm \\
10 & A batter is preparing to hit the ball at a game. & 17 $\to$ 28 & Ảnh đúng thiếu nhãn \emph{batter}, \emph{ball} \\
\bottomrule
\end{tabular}
\end{table*}

\textbf{Phép thử xoá.} Đây là mức fidelity cơ bản của cấp~1. Trên ảnh mô hình xếp hạng 1 của mỗi ca, nhóm xoá phần ảnh có
đóng góp lớn nhất vào điểm, rồi so với trung bình 10 lần xoá ngẫu nhiên một phần cùng loại.
\begin{itemize}
\item Ở 10/10 ca, xoá phần lớn nhất làm điểm giảm nhiều hơn xoá ngẫu nhiên.
\item Ở 7/10 ca, riêng việc xoá phần đó đã làm ảnh rơi khỏi hạng 1 (mô hình cấp~1: 5/10). Khi phần bị xoá là vật thể, 6/6 ca
rơi hạng; khi là bộ ba, chỉ 1/4.
\item Attention trung bình của cặp khớp nhất là 0,51 ở kênh vật thể và 0,17 ở kênh bộ ba, tức attention của kênh bộ ba
phân tán hơn.
\end{itemize}
Mục~\ref{sec:xai} đo fidelity, validity và sparsity trên toàn bộ val.

\textbf{Nguyên nhân thất bại} quan sát được khi xem ảnh thật, ở cả hai mô hình: chú thích VG thiếu nhãn hoặc đồ thị rỗng;
khác biệt từ vựng giữa caption COCO và nhãn VG (\emph{plate}/\emph{tray}, \emph{child}/\emph{boy}); thuộc tính (màu, chất
liệu) chưa được dùng; chi tiết mà scene graph không biểu diễn được (chữ trên biển hiệu); ảnh đúng nằm ngoài top-50 của CLIP.
Về thuộc tính, một phép thử không huấn luyện chỉ tăng R@1 val 0,36 điểm, dưới ngưỡng 0,5 điểm nhóm đặt trước, nên nhóm
không thêm thuộc tính vào mô hình (\file{experiments/checks/attribute-probe.REPORT.md}).
